\chapter{Limites, sécurité et suite du projet}
\label{chap:limites}

\begin{chapeau}
Ce chapitre dit ce que le logiciel ne fait pas encore, ce qu'il faut surveiller avec les données, et ce qu'il reste à
faire. Il se termine par la conclusion du projet.
\end{chapeau}

\section{Les limites}

\begin{itemize}
  \item \textbf{Un seul chien, et il est sain.} Les déformités sont simulées. Un os très déformé, un implant, un cal
  vicieux ou un petit chien pourraient tromper la détection des repères. Les repères doivent toujours être vérifiés à
  l'écran (étape 1 de l'onglet CORA, ou \fichier{cora_<os>.png}).
  \item \textbf{Les angles de profil} reposent sur des lignes articulaires courtes (20 à 30~mm) : ils sont sensibles à
  quelques degrés.
  \item \textbf{Les valeurs normales} sont publiées surtout pour quelques races (Labrador, berger allemand, golden
  retriever, rottweiler) et mesurées sur radiographie ; pour le lévrier, seule la pente du plateau tibial a une
  référence propre. Le meilleur témoin reste le membre opposé.
  \item \textbf{Un seul apex}, et au fémur seulement les déformités distales (\cref{sec:hypotheses}).
  \item \textbf{Seulement des coupes planes} : pas d'ostéotomie en dôme (\cref{sec:dome}).
  \item \textbf{Les cartilages et les physes} d'un os segmenté sont approchés par des fractions de la longueur de l'os ;
  les os voisins ne sont évités que s'ils ont été segmentés.
  \item \textbf{Les voies d'abord et les marges} (cartilage 3 à 5~mm, os voisins 1 à 2~mm, entrée des broches à
  50\degre{} au plus\ldots) sont des choix raisonnés, à valider avec un chirurgien.
  \item \textbf{Les tissus mous et la plaque} ne sont pas modélisés : une broche peut gêner la plaque, et l'abord doit
  dégager toute la zone d'appui du guide.
  \item \textbf{Les guides n'ont pas été imprimés.} Les jeux de 0,2~mm des fentes et des manchons doivent être réglés
  sur de petites pièces de test (coupons) ; la tenue à la stérilisation reste à vérifier (matériau conseillé : polyamide
  PA12 fritté au laser, ou résine stérilisable).
  \item \textbf{L'application autonome Windows} n'a pas encore été construite ; le lanceur Windows de développement,
  lui, a été essayé.
\end{itemize}

\section{La sécurité et les données}
\label{sec:securite}

\begin{itemize}
  \item \textbf{Tout reste sur l'ordinateur.} Le serveur n'écoute que \code{127.0.0.1} et l'application n'envoie rien
  sur Internet.
  \item \textbf{Les en-têtes DICOM contiennent des noms.} Le nom de l'animal, parfois celui du propriétaire, et un
  identifiant sont recopiés dans \fichier{rapport.md}, \fichier{rapport.json} et le journal. Les résultats d'exemple du
  dépôt GitHub en contiennent : il faut les rendre anonymes avant de partager le dépôt plus largement (supprimer ces
  fichiers ne suffit pas, ils restent dans l'historique Git).
  \item \textbf{Le serveur ne vérifie pas l'origine des demandes.} Une page web ouverte sur le même ordinateur pourrait,
  en théorie, lui envoyer une commande. Il faudrait un jeton tiré au hasard au démarrage et un contrôle de l'origine.
  \item \textbf{OneDrive} copie le projet sur les serveurs de l'école : c'est acceptable pour les données d'exemple ;
  pour de vrais cas cliniques, il faut l'accord de l'établissement, ou un dossier non synchronisé.
  \item \textbf{Les fichiers à imprimer} (STL) ne contiennent que la forme des pièces, aucune donnée sur l'animal.
\end{itemize}

\begin{attention}[title={Un prototype de recherche}]
CT-CORA3D n'est pas un dispositif médical. Il ne doit pas servir sur un animal sans validation de chaque planification
par le chirurgien.
\end{attention}

\section{La suite du projet}

\begin{enumerate}
  \item \textbf{Imprimer et essayer.} D'abord des coupons de fentes et de manchons pour régler les jeux, puis des
  guides sur des os imprimés ; mesurer l'erreur finale (angle, torsion, translation) par scanner de l'os corrigé.
  \item \textbf{Vrais cas.} Obtenir les scanners complets de chiens opérés, rendus anonymes (\cref{sec:topogramme}), et
  comparer nos planifications à celles du chirurgien.
  \item \textbf{Robustesse.} D'autres scanners et d'autres races ; le membre opposé comme référence quand il est
  scanné ; des os vraiment déformés.
  \item \textbf{Nouvelles fonctions.} Une plaque sur mesure dont les trous de vis sont percés à travers le guide
  \cite{carwardine2021} ; les déformités à deux apex ; l'application autonome pour Windows.
  \item \textbf{Sécurité.} Jeton et contrôle d'origine sur le serveur ; ne plus recopier les noms dans les rapports ;
  signature de l'application avant de la diffuser.
\end{enumerate}

\section{Conclusion}

CT-CORA3D relie, dans une seule application, le scanner d'une patte de chien à un guide de coupe prêt à imprimer. Chaque
étape s'appuie sur une méthode publiée ou sur un calcul simple et vérifiable : la matrice DICOM, une segmentation par
seuil montant et croissance de région, des repères trouvés comme le ferait le chirurgien avec sa règle, la méthode CORA
de Paley faite en 3D, une stratégie choisie par des règles écrites, et un guide modélisé par des champs de distance puis
découpé par des opérations exactes. Sur des cas simulés, le calcul est exact à quelques millièmes de degré, et plus de
1\,800 cas passent tous les contrôles.

Les campagnes de simulation ont été l'outil le plus utile du projet : elles ont trouvé des centaines d'erreurs que
personne n'aurait vues sur un seul cas. Mais elles ont aussi montré leur limite : un guide « valide » sur le papier peut
être impossible à poser, et c'est le regard du chirurgien qui l'a révélé. La prochaine étape n'est donc plus seulement
informatique : il faut imprimer, essayer, et comparer avec le chirurgien sur de vrais patients.
